% Options for packages loaded elsewhere
\PassOptionsToPackage{unicode}{hyperref}
\PassOptionsToPackage{hyphens}{url}
\documentclass[
  a4paper,
]{ltjsarticle}
\usepackage{xcolor}
\usepackage[margin=25mm]{geometry}
\usepackage{amsmath,amssymb}
\setcounter{secnumdepth}{-\maxdimen} % remove section numbering
\usepackage{iftex}
\ifPDFTeX
  \usepackage[T1]{fontenc}
  \usepackage[utf8]{inputenc}
  \usepackage{textcomp} % provide euro and other symbols
\else % if luatex or xetex
  \usepackage{unicode-math} % this also loads fontspec
  \defaultfontfeatures{Scale=MatchLowercase}
  \defaultfontfeatures[\rmfamily]{Ligatures=TeX,Scale=1}
\fi
\usepackage{lmodern}
\ifPDFTeX\else
  % xetex/luatex font selection
\fi
% Use upquote if available, for straight quotes in verbatim environments
\IfFileExists{upquote.sty}{\usepackage{upquote}}{}
\IfFileExists{microtype.sty}{% use microtype if available
  \usepackage[]{microtype}
  \UseMicrotypeSet[protrusion]{basicmath} % disable protrusion for tt fonts
}{}
\makeatletter
\@ifundefined{KOMAClassName}{% if non-KOMA class
  \IfFileExists{parskip.sty}{%
    \usepackage{parskip}
  }{% else
    \setlength{\parindent}{0pt}
    \setlength{\parskip}{6pt plus 2pt minus 1pt}}
}{% if KOMA class
  \KOMAoptions{parskip=half}}
\makeatother
\usepackage{color}
\usepackage{fancyvrb}
\newcommand{\VerbBar}{|}
\newcommand{\VERB}{\Verb[commandchars=\\\{\}]}
\DefineVerbatimEnvironment{Highlighting}{Verbatim}{commandchars=\\\{\}}
% Add ',fontsize=\small' for more characters per line
\newenvironment{Shaded}{}{}
\newcommand{\AlertTok}[1]{\textcolor[rgb]{1.00,0.00,0.00}{\textbf{#1}}}
\newcommand{\AnnotationTok}[1]{\textcolor[rgb]{0.38,0.63,0.69}{\textbf{\textit{#1}}}}
\newcommand{\AttributeTok}[1]{\textcolor[rgb]{0.49,0.56,0.16}{#1}}
\newcommand{\BaseNTok}[1]{\textcolor[rgb]{0.25,0.63,0.44}{#1}}
\newcommand{\BuiltInTok}[1]{\textcolor[rgb]{0.00,0.50,0.00}{#1}}
\newcommand{\CharTok}[1]{\textcolor[rgb]{0.25,0.44,0.63}{#1}}
\newcommand{\CommentTok}[1]{\textcolor[rgb]{0.38,0.63,0.69}{\textit{#1}}}
\newcommand{\CommentVarTok}[1]{\textcolor[rgb]{0.38,0.63,0.69}{\textbf{\textit{#1}}}}
\newcommand{\ConstantTok}[1]{\textcolor[rgb]{0.53,0.00,0.00}{#1}}
\newcommand{\ControlFlowTok}[1]{\textcolor[rgb]{0.00,0.44,0.13}{\textbf{#1}}}
\newcommand{\DataTypeTok}[1]{\textcolor[rgb]{0.56,0.13,0.00}{#1}}
\newcommand{\DecValTok}[1]{\textcolor[rgb]{0.25,0.63,0.44}{#1}}
\newcommand{\DocumentationTok}[1]{\textcolor[rgb]{0.73,0.13,0.13}{\textit{#1}}}
\newcommand{\ErrorTok}[1]{\textcolor[rgb]{1.00,0.00,0.00}{\textbf{#1}}}
\newcommand{\ExtensionTok}[1]{#1}
\newcommand{\FloatTok}[1]{\textcolor[rgb]{0.25,0.63,0.44}{#1}}
\newcommand{\FunctionTok}[1]{\textcolor[rgb]{0.02,0.16,0.49}{#1}}
\newcommand{\ImportTok}[1]{\textcolor[rgb]{0.00,0.50,0.00}{\textbf{#1}}}
\newcommand{\InformationTok}[1]{\textcolor[rgb]{0.38,0.63,0.69}{\textbf{\textit{#1}}}}
\newcommand{\KeywordTok}[1]{\textcolor[rgb]{0.00,0.44,0.13}{\textbf{#1}}}
\newcommand{\NormalTok}[1]{#1}
\newcommand{\OperatorTok}[1]{\textcolor[rgb]{0.40,0.40,0.40}{#1}}
\newcommand{\OtherTok}[1]{\textcolor[rgb]{0.00,0.44,0.13}{#1}}
\newcommand{\PreprocessorTok}[1]{\textcolor[rgb]{0.74,0.48,0.00}{#1}}
\newcommand{\RegionMarkerTok}[1]{#1}
\newcommand{\SpecialCharTok}[1]{\textcolor[rgb]{0.25,0.44,0.63}{#1}}
\newcommand{\SpecialStringTok}[1]{\textcolor[rgb]{0.73,0.40,0.53}{#1}}
\newcommand{\StringTok}[1]{\textcolor[rgb]{0.25,0.44,0.63}{#1}}
\newcommand{\VariableTok}[1]{\textcolor[rgb]{0.10,0.09,0.49}{#1}}
\newcommand{\VerbatimStringTok}[1]{\textcolor[rgb]{0.25,0.44,0.63}{#1}}
\newcommand{\WarningTok}[1]{\textcolor[rgb]{0.38,0.63,0.69}{\textbf{\textit{#1}}}}
\usepackage{longtable,booktabs,array}
\newcounter{none} % for unnumbered tables
\usepackage{calc} % for calculating minipage widths
% Correct order of tables after \paragraph or \subparagraph
\usepackage{etoolbox}
\makeatletter
\patchcmd\longtable{\par}{\if@noskipsec\mbox{}\fi\par}{}{}
\makeatother
% Allow footnotes in longtable head/foot
\IfFileExists{footnotehyper.sty}{\usepackage{footnotehyper}}{\usepackage{footnote}}
\makesavenoteenv{longtable}
\setlength{\emergencystretch}{3em} % prevent overfull lines
\providecommand{\tightlist}{%
  \setlength{\itemsep}{0pt}\setlength{\parskip}{0pt}}
\usepackage{bookmark}
\IfFileExists{xurl.sty}{\usepackage{xurl}}{} % add URL line breaks if available
\urlstyle{same}
\hypersetup{
  pdftitle={Chapter 1: Generation of the Static Parent Data --- Self-Consistent Circularly Polarized Eigenmode Parents and Zero-Closure Kernel Seeds for N=3..40},
  pdfauthor={Noriaki Kihara (WF System Co., Ltd.) ORCID 0009-0004-6753-4020},
  hidelinks,
  pdfcreator={LaTeX via pandoc}}

\title{Chapter 1: Generation of the Static Parent Data ---
Self-Consistent Circularly Polarized Eigenmode Parents and Zero-Closure
Kernel Seeds for N=3..40}
\author{Noriaki Kihara (WF System Co., Ltd.) ORCID 0009-0004-6753-4020}
\date{September 5, 2026 Version DOI 10.5281/zenodo.22317636}

\begin{document}
\maketitle

(Chapter 1 of the reproduction papers for the N=3..40 stage-1+2+3 sweep.
Shares the Concept DOI 10.5281/zenodo.22317635 with the overview paper;
Version DOI 10.5281/zenodo.22317636. Code blocks are quoted verbatim
from the programs; comments inside them remain in Japanese, as in the
originals.)

\subsection{1. Purpose}\label{purpose}

This chapter describes the numerical experiment that deterministically
generates, from a random-seed formula, the initial data used by the
sweep of Chapter 2 --- the self-consistent circularly polarized
eigenmode parent v, the zero-closure kernel seed g, and the normalized
initial state Z0 --- and fixes them as static files (npz).

The generation procedure, the order of random-number consumption, and
all arguments are identical to the canonical run of 2026-07-22
(\texttt{自発的分裂予備実験\_v1/run\_spontaneous\_splitting\_largeN\_v1.py};
N=40, δ=1e-15, seed=0, tol=1e-12), and the pass gate is that, for N=40,
the generated objects are bit-identical to the initial values of the
canonical run. The purpose of this chapter is complete reproducibility
and the fixing of the correspondence among equations, programs, and
data; it contains no physical interpretation.

\subsection{2. Theoretical Background}\label{theoretical-background}

We work with complex relation waves on the edge set of the complete
graph K\_N.

\textbf{(Eq. 1) Edge set and dimension} --- The number of edges is M =
N(N−1)/2. The edge ordering is fixed to the upper-triangular order
((0,1),(0,2),\ldots,(0,N−1),(1,2),\ldots). The state is z ∈ C\^{}M.

\textbf{(Eq. 2) Phase-difference sine generator} --- For edge phases θ ∈
R\^{}M, on vertex-sharing edge pairs (e,f):

\begin{verbatim}
K_ef(θ) = cos θ_e sin θ_f − sin θ_e cos θ_f = sin(θ_f − θ_e)
\end{verbatim}

0 on non-adjacent pairs, 0 on the diagonal. K is real antisymmetric
(K\^{}T = −K).

\textbf{(Eq. 3) Vertex decomposition and the small space} --- Scattering
c\_e = cos θ\_e, s\_e = sin θ\_e to the vertices yields matrices C, S
(M×N); with W = {[}C\textbar S{]} (M×2N) and J = {[}{[}0, I\_N{]},
{[}−I\_N, 0{]}{]}:

\begin{verbatim}
K = C S^T − S C^T = W J W^T,   rank K ≤ 2N
\end{verbatim}

The Gram matrix G = W\^{}T W (2N×2N) can be constructed analytically
because, in the line graph of K\_N, a vertex pair (k,l) shares exactly
one edge (the edge (k,l)):

\begin{verbatim}
G_cc[k,l] = cos^2θ_{(k,l)},  G_cs[k,l] = cosθ_{(k,l)} sinθ_{(k,l)},  G_ss[k,l] = sin^2θ_{(k,l)}  (k≠l)
diagonals are the row sums of the off-diagonals
\end{verbatim}

\textbf{(Eq. 4) Circularly polarized eigenmode (parent)} --- The nonzero
spectrum of K corresponds to the eigenvalues of JG (2N×2N). From the
eigenvector y of the eigenvalue λ with the smallest imaginary part (λ =
−iσ\_max):

\begin{verbatim}
v = W y,   v ← v/‖v‖
\end{verbatim}

Since iK is Hermitian, if v is an eigenmode then iKv = μv (μ = σ\_max).

\textbf{(Eq. 5) Self-consistent iteration (phase mixing)} --- The phases
θ\_new = arg v read from v are mixed into the current phases with mixing
ratio β = 0.5:

\begin{verbatim}
θ ← arg( (1−β) e^{iθ} + β e^{iθ_new} )
\end{verbatim}

A fixed point of this iteration is the self-consistent parent: ``the
eigenmode of the generator built from its own phases is itself''.

\textbf{(Eq. 6) Eigenmode residual (convergence criterion)} ---

\begin{verbatim}
μ = Re( v† (iKv) ),   r = ‖ iKv − μ v ‖
\end{verbatim}

Convergence is declared when r \textless{} tol (tol = 1e-12 in this
chapter).

\textbf{(Eq. 7) Zero-closure kernel seed} --- With the projection onto
the orthogonal complement of span(W)

\begin{verbatim}
P⊥ g = g − W q,   q = lstsq(G, W^T g)
\end{verbatim}

construct from real Gaussian vectors ξ\_1, ξ\_2: u = P⊥ξ\_1/‖·‖ and w =
the orthonormalized P⊥ξ\_2, and set

\begin{verbatim}
g = (u + i w)/√2
\end{verbatim}

From u ⊥ w and ‖u‖ = ‖w‖ = 1, g\^{}Tg = (‖u‖\^{}2 − ‖w‖\^{}2)/2 + i u·w
= 0 holds exactly (zero closure). Moreover g ∈ span(W)⊥, i.e., g lies
outside the range of K.

\textbf{(Eq. 8) Initial state} ---

\begin{verbatim}
Z0 = (v + δ g) / ‖v + δ g‖,   δ = 1e-15
\end{verbatim}

\textbf{(Eq. 9) Random numbers} --- The RNG is numpy PCG64 with seed

\begin{verbatim}
seed(N) = 40260722 + 1000·N + 0
\end{verbatim}

The consumption order is fixed: ``initial phases θ\^{}0 of the parent
iteration (m uniform numbers, per restart) → ξ\_1 (m normal numbers) →
ξ\_2 (m normal numbers)''. All generated objects are thereby
deterministic functions of N alone.

\subsection{3. Implementation Method}\label{implementation-method}

\begin{itemize}
\tightlist
\item
  The engine \texttt{run\_n\_scaling\_lowrank\_v1.py} is a
  \textbf{bit-identical copy} of the 2026-07-22 canonical
  (\texttt{自発的分裂予備実験\_v1/run\_n\_scaling\_lowrank\_v1.py}),
  verified by diff. The implementations of Eqs. 2--7 all reside in this
  engine; this chapter uses them \textbf{unmodified, via import only}
  (following the rule prohibiting independent reimplementation).
\item
  The generator program \texttt{make\_static\_parents\_N3\_N40\_v1.py}
  executes, for N=3..40, the same call sequence as the opening of the
  canonical \texttt{run()} (rng of Eq. 9 → make\_parent →
  zero\_closure\_kernel\_seed → normalization of Eq. 8) and saves to
  npz.
\item
  For N=40, bit-identity with the canonical static parent (2026-09-04;
  itself verified bit-identical against the canonical run) is checked as
  an in-program gate.
\end{itemize}

\subsection{4. Detailed Design}\label{detailed-design}

\subsubsection{4.1 Overall Flow}\label{overall-flow}

\begin{verbatim}
for N in 3..40:
  [initialization]  build LowRankSystem(N) (edge order, J), rng = default_rng(40260722+1000N)
  [loop]  make_parent: self-consistent iteration (Eq. 4 → Eq. 5 → Eq. 6; up to 1200 iterations × 3 restarts)
          zero_closure_kernel_seed: construction of the seed g (Eq. 7)
  [finalization]  build Z0 (Eq. 8) → N=40 gate → save npz (verify only if file exists) → append to the ledger
write the ledger parents_summary.csv; exit 1 if any gate fails
\end{verbatim}

\subsubsection{4.2 Overall Data Flow}\label{overall-data-flow}

\begin{itemize}
\tightlist
\item
  \textbf{Input}: no file input. Only the random-seed formula (Eq. 9)
  and constants.

  \begin{itemize}
  \tightlist
  \item
    \texttt{SEED\ =\ 0} \ldots{} third term of the seed formula (series
    number)
  \item
    \texttt{DELTA\ =\ 1e-15} \ldots{} δ of Eq. 8 (seed amplitude)
  \item
    \texttt{TOL\ =\ 1e-12} \ldots{} convergence threshold of Eq. 6
    (identical to the canonical run's \texttt{-\/-tol=1e-12})
  \item
    \texttt{ITERS\ =\ 1200} \ldots{} iteration cap of the
    self-consistent loop (identical to the canonical \texttt{run()}'s
    \texttt{iters=1200})
  \item
    reference input \texttt{REF40} \ldots{} the canonical static-parent
    npz for the N=40 gate (read-only)
  \end{itemize}
\item
  \textbf{Processing}: repeat the individual processes of 4.3 for
  N=3..40.
\item
  \textbf{Output}:

  \begin{itemize}
  \tightlist
  \item
    \texttt{parents/parent\_static\_N\{N:05d\}\_makeparent\_20260905.npz}
    × 38 (fields: \texttt{v}, \texttt{g}, \texttt{Z0}, \texttt{sigma}
    (the parent's σ spectrum), \texttt{residual}, \texttt{n},
    \texttt{seed}, \texttt{delta}, \texttt{tol}, \texttt{iters})
  \item
    \texttt{parents/parents\_summary.csv} (columns: N, M,
    parent\_residual, rank\_planes, status)
  \end{itemize}
\end{itemize}

Constant definitions (\texttt{make\_static\_parents\_N3\_N40\_v1.py}):

\begin{Shaded}
\begin{Highlighting}[]
    \DecValTok{34}\NormalTok{  PARENT\_DIR }\OperatorTok{=}\NormalTok{ os.path.join(BASE\_DIR, }\StringTok{"parents"}\NormalTok{)}
    \DecValTok{35}\NormalTok{  REF40 }\OperatorTok{=} \StringTok{\textquotesingle{}.../自発的分裂予備実験\_v1\_N40対照実験系\_20260904/largeN\_splitting\_result\_v1/parent\_static\_N40\_makeparent\_20260904.npz\textquotesingle{}}
    \DecValTok{36}
    \DecValTok{37}\NormalTok{  SEED }\OperatorTok{=} \DecValTok{0}
    \DecValTok{38}\NormalTok{  DELTA }\OperatorTok{=} \FloatTok{1e{-}15}
    \DecValTok{39}\NormalTok{  TOL }\OperatorTok{=} \FloatTok{1e{-}12}
    \DecValTok{40}\NormalTok{  ITERS }\OperatorTok{=} \DecValTok{1200}
\end{Highlighting}
\end{Shaded}

\subsubsection{4.3 Individual Processes}\label{individual-processes}

\paragraph{4.3.1 Initialization}\label{initialization}

Engine initialization (Eq. 1; J of Eq. 3). The edge order is fixed by
\texttt{np.triu\_indices}.

\texttt{run\_n\_scaling\_lowrank\_v1.py}:

\begin{Shaded}
\begin{Highlighting}[]
    \DecValTok{52}  \KeywordTok{def}\NormalTok{ build\_edges(n):}
    \DecValTok{53}\NormalTok{      ea, eb }\OperatorTok{=}\NormalTok{ np.triu\_indices(n, k}\OperatorTok{=}\DecValTok{1}\NormalTok{)}
    \DecValTok{54}      \ControlFlowTok{return}\NormalTok{ ea.astype(np.int64), eb.astype(np.int64)}
\NormalTok{...}
    \DecValTok{60}      \KeywordTok{def} \FunctionTok{\_\_init\_\_}\NormalTok{(}\VariableTok{self}\NormalTok{, n):}
    \DecValTok{61}          \VariableTok{self}\NormalTok{.n }\OperatorTok{=}\NormalTok{ n}
    \DecValTok{62}          \VariableTok{self}\NormalTok{.ea, }\VariableTok{self}\NormalTok{.eb }\OperatorTok{=}\NormalTok{ build\_edges(n)}
    \DecValTok{63}          \VariableTok{self}\NormalTok{.m }\OperatorTok{=} \BuiltInTok{len}\NormalTok{(}\VariableTok{self}\NormalTok{.ea)}
    \DecValTok{64}          \VariableTok{self}\NormalTok{.J }\OperatorTok{=}\NormalTok{ np.zeros((}\DecValTok{2} \OperatorTok{*}\NormalTok{ n, }\DecValTok{2} \OperatorTok{*}\NormalTok{ n))}
    \DecValTok{65}          \VariableTok{self}\NormalTok{.J[:n, n:] }\OperatorTok{=}\NormalTok{ np.eye(n)}
    \DecValTok{66}          \VariableTok{self}\NormalTok{.J[n:, :n] }\OperatorTok{=} \OperatorTok{{-}}\NormalTok{np.eye(n)}
\end{Highlighting}
\end{Shaded}

Caller side (\texttt{make\_static\_parents\_N3\_N40\_v1.py}, Eq. 9):

\begin{Shaded}
\begin{Highlighting}[]
    \DecValTok{48}\NormalTok{          sys\_lr }\OperatorTok{=}\NormalTok{ LowRankSystem(N)}
    \DecValTok{49}\NormalTok{          rng }\OperatorTok{=}\NormalTok{ np.random.default\_rng(}\DecValTok{40260722} \OperatorTok{+} \DecValTok{1000} \OperatorTok{*}\NormalTok{ N }\OperatorTok{+}\NormalTok{ SEED)}
\end{Highlighting}
\end{Shaded}

\paragraph{4.3.2 Loop (1): Construction of the Self-Consistent Parent
(Eqs.
2--6)}\label{loop-1-construction-of-the-self-consistent-parent-eqs.-26}

Generator construction (Eqs. 2, 3; \texttt{set\_theta} holds c, s and
the analytic G):

\begin{Shaded}
\begin{Highlighting}[]
    \DecValTok{68}      \KeywordTok{def}\NormalTok{ set\_theta(}\VariableTok{self}\NormalTok{, theta):}
    \DecValTok{69}\NormalTok{          n }\OperatorTok{=} \VariableTok{self}\NormalTok{.n}
    \DecValTok{70}          \VariableTok{self}\NormalTok{.c }\OperatorTok{=}\NormalTok{ np.cos(theta)}
    \DecValTok{71}          \VariableTok{self}\NormalTok{.s }\OperatorTok{=}\NormalTok{ np.sin(theta)}
    \DecValTok{72}\NormalTok{          T }\OperatorTok{=}\NormalTok{ np.zeros((n, n))}
    \DecValTok{73}\NormalTok{          T[}\VariableTok{self}\NormalTok{.ea, }\VariableTok{self}\NormalTok{.eb] }\OperatorTok{=}\NormalTok{ theta}
    \DecValTok{74}\NormalTok{          T[}\VariableTok{self}\NormalTok{.eb, }\VariableTok{self}\NormalTok{.ea] }\OperatorTok{=}\NormalTok{ theta}
    \DecValTok{75}\NormalTok{          CT }\OperatorTok{=}\NormalTok{ np.cos(T)}
    \DecValTok{76}\NormalTok{          ST }\OperatorTok{=}\NormalTok{ np.sin(T)}
    \DecValTok{77}\NormalTok{          np.fill\_diagonal(CT, }\FloatTok{0.0}\NormalTok{)}
    \DecValTok{78}\NormalTok{          np.fill\_diagonal(ST, }\FloatTok{0.0}\NormalTok{)}
    \DecValTok{79}\NormalTok{          Gcc }\OperatorTok{=}\NormalTok{ CT }\OperatorTok{*}\NormalTok{ CT}
    \DecValTok{80}\NormalTok{          Gcs }\OperatorTok{=}\NormalTok{ CT }\OperatorTok{*}\NormalTok{ ST}
    \DecValTok{81}\NormalTok{          Gss }\OperatorTok{=}\NormalTok{ ST }\OperatorTok{*}\NormalTok{ ST}
    \DecValTok{82}\NormalTok{          np.fill\_diagonal(Gcc, Gcc.}\BuiltInTok{sum}\NormalTok{(axis}\OperatorTok{=}\DecValTok{1}\NormalTok{))}
    \DecValTok{83}\NormalTok{          np.fill\_diagonal(Gcs, Gcs.}\BuiltInTok{sum}\NormalTok{(axis}\OperatorTok{=}\DecValTok{1}\NormalTok{))}
    \DecValTok{84}\NormalTok{          np.fill\_diagonal(Gss, Gss.}\BuiltInTok{sum}\NormalTok{(axis}\OperatorTok{=}\DecValTok{1}\NormalTok{))}
    \DecValTok{85}\NormalTok{          G }\OperatorTok{=}\NormalTok{ np.empty((}\DecValTok{2} \OperatorTok{*}\NormalTok{ n, }\DecValTok{2} \OperatorTok{*}\NormalTok{ n))}
    \DecValTok{86}\NormalTok{          G[:n, :n] }\OperatorTok{=}\NormalTok{ Gcc}
    \DecValTok{87}\NormalTok{          G[:n, n:] }\OperatorTok{=}\NormalTok{ Gcs}
    \DecValTok{88}\NormalTok{          G[n:, :n] }\OperatorTok{=}\NormalTok{ Gcs}
    \DecValTok{89}\NormalTok{          G[n:, n:] }\OperatorTok{=}\NormalTok{ Gss}
    \DecValTok{90}          \VariableTok{self}\NormalTok{.G }\OperatorTok{=}\NormalTok{ G}
\end{Highlighting}
\end{Shaded}

Parent iteration body (Eq. 4: lines 169--172; Eq. 5: lines 173--175; the
criterion of Eq. 6: lines 176--181):

\begin{Shaded}
\begin{Highlighting}[]
   \DecValTok{158}  \KeywordTok{def}\NormalTok{ make\_parent(sys\_lr, rng, iters}\OperatorTok{=}\DecValTok{400}\NormalTok{, beta}\OperatorTok{=}\FloatTok{0.5}\NormalTok{, tol}\OperatorTok{=}\FloatTok{1e{-}8}\NormalTok{, restarts}\OperatorTok{=}\DecValTok{3}\NormalTok{):}
   \DecValTok{159}      \StringTok{"""自己無撞着円偏波固有モード親。小空間（JG）の固有対で反復。}
\StringTok{   160}
\StringTok{   161      収束判定（残差 \textless{} tol）付き。停滞時はランダム初期位相からリスタート。}
\StringTok{   162      """}
   \DecValTok{163}\NormalTok{      best }\OperatorTok{=}\NormalTok{ (}\VariableTok{None}\NormalTok{, np.inf, }\VariableTok{None}\NormalTok{)}
   \DecValTok{164}      \ControlFlowTok{for}\NormalTok{ \_ }\KeywordTok{in} \BuiltInTok{range}\NormalTok{(restarts):}
   \DecValTok{165}\NormalTok{          theta }\OperatorTok{=}\NormalTok{ rng.uniform(}\FloatTok{0.0}\NormalTok{, }\FloatTok{2.0} \OperatorTok{*}\NormalTok{ np.pi, sys\_lr.m)}
   \DecValTok{166}\NormalTok{          v }\OperatorTok{=} \VariableTok{None}
   \DecValTok{167}          \ControlFlowTok{for}\NormalTok{ it }\KeywordTok{in} \BuiltInTok{range}\NormalTok{(iters):}
   \DecValTok{168}\NormalTok{              sys\_lr.set\_theta(theta)}
   \DecValTok{169}\NormalTok{              ev, EV }\OperatorTok{=}\NormalTok{ np.linalg.eig(sys\_lr.J }\OperatorTok{@}\NormalTok{ sys\_lr.G)}
   \DecValTok{170}\NormalTok{              idx }\OperatorTok{=} \BuiltInTok{int}\NormalTok{(np.argmin(ev.imag))  }\CommentTok{\# λ = {-}iσ\_max}
   \DecValTok{171}\NormalTok{              v }\OperatorTok{=}\NormalTok{ sys\_lr.w(EV[:, idx].astype(}\BuiltInTok{complex}\NormalTok{))}
   \DecValTok{172}\NormalTok{              v }\OperatorTok{=}\NormalTok{ v }\OperatorTok{/}\NormalTok{ np.linalg.norm(v)}
   \DecValTok{173}\NormalTok{              theta\_new }\OperatorTok{=}\NormalTok{ np.angle(v)}
   \DecValTok{174}\NormalTok{              mix }\OperatorTok{=}\NormalTok{ (}\FloatTok{1.0} \OperatorTok{{-}}\NormalTok{ beta) }\OperatorTok{*}\NormalTok{ np.exp(}\OtherTok{1j} \OperatorTok{*}\NormalTok{ theta) }\OperatorTok{+}\NormalTok{ beta }\OperatorTok{*}\NormalTok{ np.exp(}\OtherTok{1j} \OperatorTok{*}\NormalTok{ theta\_new)}
   \DecValTok{175}\NormalTok{              theta }\OperatorTok{=}\NormalTok{ np.angle(mix)}
   \DecValTok{176}              \ControlFlowTok{if}\NormalTok{ it }\OperatorTok{\%} \DecValTok{10} \OperatorTok{==} \DecValTok{9}\NormalTok{:}
   \DecValTok{177}\NormalTok{                  sys\_lr.set\_theta(np.angle(v))}
   \DecValTok{178}\NormalTok{                  res\_now }\OperatorTok{=}\NormalTok{ \_eigenmode\_residual(sys\_lr, v)}
   \DecValTok{179}\NormalTok{                  progress(}\SpecialStringTok{f"親構成 iter=}\SpecialCharTok{\{}\NormalTok{it}\OperatorTok{+}\DecValTok{1}\SpecialCharTok{\}}\SpecialStringTok{ 残差=}\SpecialCharTok{\{}\NormalTok{res\_now}\SpecialCharTok{:.2e\}}\SpecialStringTok{"}\NormalTok{)}
   \DecValTok{180}                  \ControlFlowTok{if}\NormalTok{ res\_now }\OperatorTok{\textless{}}\NormalTok{ tol:}
   \DecValTok{181}                      \ControlFlowTok{break}
   \DecValTok{182}\NormalTok{          sys\_lr.set\_theta(np.angle(v))}
   \DecValTok{183}\NormalTok{          residual }\OperatorTok{=}\NormalTok{ \_eigenmode\_residual(sys\_lr, v)}
   \DecValTok{184}          \ControlFlowTok{if}\NormalTok{ residual }\OperatorTok{\textless{}}\NormalTok{ best[}\DecValTok{1}\NormalTok{]:}
   \DecValTok{185}\NormalTok{              best }\OperatorTok{=}\NormalTok{ (v, residual, sys\_lr.sigma\_spectrum())}
   \DecValTok{186}          \ControlFlowTok{if}\NormalTok{ residual }\OperatorTok{\textless{}}\NormalTok{ tol:}
   \DecValTok{187}              \ControlFlowTok{break}
   \DecValTok{188}\NormalTok{      v, residual, sig }\OperatorTok{=}\NormalTok{ best}
   \DecValTok{189}\NormalTok{      sys\_lr.set\_theta(np.angle(v))}
   \DecValTok{190}      \ControlFlowTok{return}\NormalTok{ v, residual, sig}
\end{Highlighting}
\end{Shaded}

Residual (Eq. 6):

\begin{Shaded}
\begin{Highlighting}[]
   \DecValTok{151}  \KeywordTok{def}\NormalTok{ \_eigenmode\_residual(sys\_lr, v):}
   \DecValTok{152}      \StringTok{"""カイラリティ非依存の固有モード残差: μ = v†(iKv) に対する ‖iKv {-} μv‖。"""}
   \DecValTok{153}\NormalTok{      kv }\OperatorTok{=}\NormalTok{ sys\_lr.kmatvec(v)}
   \DecValTok{154}\NormalTok{      mu }\OperatorTok{=} \BuiltInTok{float}\NormalTok{(np.real(np.conj(v) }\OperatorTok{@}\NormalTok{ (}\OtherTok{1j} \OperatorTok{*}\NormalTok{ kv)))}
   \DecValTok{155}      \ControlFlowTok{return} \BuiltInTok{float}\NormalTok{(np.linalg.norm(}\OtherTok{1j} \OperatorTok{*}\NormalTok{ kv }\OperatorTok{{-}}\NormalTok{ mu }\OperatorTok{*}\NormalTok{ v))}
\end{Highlighting}
\end{Shaded}

\textbf{Stopping conditions and exceptional behavior (elements not
formulated as equations)}: - The convergence check is performed only
\textbf{every 10 iterations} (line 176, \texttt{it\ \%\ 10\ ==\ 9}); the
actual iteration count therefore ends at a multiple of 10. - If r
\textless{} tol is not reached within \texttt{iters=1200}, no exception
is raised; the loop proceeds to a restart (line 165, consuming new
initial phases from the rng), up to \texttt{restarts=3}. The solution
with the best residual is adopted (lines 163, 184--185). \textbf{A
restart changes the number of rng draws}, so reproduction requires the
same tol. (In this run, every N converged within the first restart with
no extra draws; evidence: the N=40 products are bit-identical to the
canonical generated under the same conditions.) - The \texttt{progress}
line (line 179) is stderr output only and does not affect the
computation.

\paragraph{4.3.3 Loop (2): Construction of the Zero-Closure Kernel Seed
(Eq.
7)}\label{loop-2-construction-of-the-zero-closure-kernel-seed-eq.-7}

\begin{Shaded}
\begin{Highlighting}[]
   \DecValTok{193}  \KeywordTok{def}\NormalTok{ zero\_closure\_kernel\_seed(sys\_lr, rng):}
   \DecValTok{194}      \StringTok{"""span(W) の直交補内の実正規直交対 (u,w) から g=(u+iw)/√2。g\^{}T g = 0 厳密。"""}
   \DecValTok{195}\NormalTok{      m }\OperatorTok{=}\NormalTok{ sys\_lr.m}
   \DecValTok{196}      \KeywordTok{def}\NormalTok{ project\_out(g):}
   \DecValTok{197}\NormalTok{          q }\OperatorTok{=}\NormalTok{ np.linalg.lstsq(sys\_lr.G, sys\_lr.wt(g), rcond}\OperatorTok{=}\VariableTok{None}\NormalTok{)[}\DecValTok{0}\NormalTok{]}
   \DecValTok{198}          \ControlFlowTok{return}\NormalTok{ g }\OperatorTok{{-}}\NormalTok{ sys\_lr.w(q)}
   \DecValTok{199}\NormalTok{      u }\OperatorTok{=}\NormalTok{ project\_out(rng.normal(size}\OperatorTok{=}\NormalTok{m))}
   \DecValTok{200}\NormalTok{      u }\OperatorTok{=}\NormalTok{ u }\OperatorTok{/}\NormalTok{ np.linalg.norm(u)}
   \DecValTok{201}\NormalTok{      w }\OperatorTok{=}\NormalTok{ project\_out(rng.normal(size}\OperatorTok{=}\NormalTok{m))}
   \DecValTok{202}\NormalTok{      w }\OperatorTok{=}\NormalTok{ w }\OperatorTok{{-}}\NormalTok{ (w }\OperatorTok{@}\NormalTok{ u) }\OperatorTok{*}\NormalTok{ u}
   \DecValTok{203}\NormalTok{      w }\OperatorTok{=}\NormalTok{ w }\OperatorTok{/}\NormalTok{ np.linalg.norm(w)}
   \DecValTok{204}      \ControlFlowTok{return}\NormalTok{ (u }\OperatorTok{+} \OtherTok{1j} \OperatorTok{*}\NormalTok{ w) }\OperatorTok{/}\NormalTok{ math.sqrt(}\FloatTok{2.0}\NormalTok{)}
\end{Highlighting}
\end{Shaded}

\begin{itemize}
\tightlist
\item
  The projection (P⊥ of Eq. 7) solves the normal equations G q = W\^{}T
  g by \texttt{lstsq} (rcond=None; line 197). Even for near-singular G,
  lstsq returns a least-squares solution, so no exception arises.
\item
  The G at this point is that of \textbf{θ = arg v} set at the end of
  make\_parent (line 189); i.e., the seed is taken from the orthogonal
  complement of the range of the parent's generator.
\end{itemize}

\paragraph{4.3.4 Finalization: Z0 Construction, Gate, Saving,
Ledger}\label{finalization-z0-construction-gate-saving-ledger}

\texttt{make\_static\_parents\_N3\_N40\_v1.py} (Eq. 8: lines 52--53):

\begin{Shaded}
\begin{Highlighting}[]
    \DecValTok{50}\NormalTok{          v, residual, sig }\OperatorTok{=}\NormalTok{ make\_parent(sys\_lr, rng, iters}\OperatorTok{=}\NormalTok{ITERS, tol}\OperatorTok{=}\NormalTok{TOL)}
    \DecValTok{51}\NormalTok{          g }\OperatorTok{=}\NormalTok{ zero\_closure\_kernel\_seed(sys\_lr, rng)}
    \DecValTok{52}\NormalTok{          Z }\OperatorTok{=}\NormalTok{ v }\OperatorTok{+}\NormalTok{ DELTA }\OperatorTok{*}\NormalTok{ g}
    \DecValTok{53}\NormalTok{          Z }\OperatorTok{=}\NormalTok{ Z }\OperatorTok{/}\NormalTok{ np.linalg.norm(Z)}
    \DecValTok{54}\NormalTok{          converged }\OperatorTok{=} \BuiltInTok{bool}\NormalTok{(residual }\OperatorTok{\textless{}} \FloatTok{1e{-}8}\NormalTok{)}
    \DecValTok{55}\NormalTok{          status }\OperatorTok{=} \StringTok{"ok"} \ControlFlowTok{if}\NormalTok{ converged }\ControlFlowTok{else} \StringTok{"NOT\_CONVERGED"}
    \DecValTok{56}
    \DecValTok{57}          \ControlFlowTok{if}\NormalTok{ N }\OperatorTok{==} \DecValTok{40}\NormalTok{:}
    \DecValTok{58}\NormalTok{              ref }\OperatorTok{=}\NormalTok{ np.load(REF40)}
    \DecValTok{59}\NormalTok{              same }\OperatorTok{=} \BuiltInTok{all}\NormalTok{(np.array\_equal(x, ref[k]) }\ControlFlowTok{for}\NormalTok{ x, k }\KeywordTok{in}\NormalTok{ ((v, }\StringTok{\textquotesingle{}v\textquotesingle{}}\NormalTok{), (g, }\StringTok{\textquotesingle{}g\textquotesingle{}}\NormalTok{), (Z, }\StringTok{\textquotesingle{}Z0\textquotesingle{}}\NormalTok{)))}
    \DecValTok{60}              \BuiltInTok{print}\NormalTok{(}\SpecialStringTok{f"GATE N=40 v/g/Z0 bit{-}identical to canonical static parent: }\SpecialCharTok{\{}\NormalTok{same}\SpecialCharTok{\}}\SpecialStringTok{"}\NormalTok{)}
    \DecValTok{61}              \ControlFlowTok{if} \KeywordTok{not}\NormalTok{ same:}
    \DecValTok{62}\NormalTok{                  gate\_ok }\OperatorTok{=} \VariableTok{False}
    \DecValTok{63}\NormalTok{                  status }\OperatorTok{=} \StringTok{"GATE\_FAIL"}
    \DecValTok{64}
    \DecValTok{65}          \ControlFlowTok{if}\NormalTok{ os.path.exists(out\_path):}
    \DecValTok{66}\NormalTok{              prev }\OperatorTok{=}\NormalTok{ np.load(out\_path)}
    \DecValTok{67}\NormalTok{              same\_prev }\OperatorTok{=} \BuiltInTok{all}\NormalTok{(np.array\_equal(x, prev[k]) }\ControlFlowTok{for}\NormalTok{ x, k }\KeywordTok{in}\NormalTok{ ((v, }\StringTok{\textquotesingle{}v\textquotesingle{}}\NormalTok{), (g, }\StringTok{\textquotesingle{}g\textquotesingle{}}\NormalTok{), (Z, }\StringTok{\textquotesingle{}Z0\textquotesingle{}}\NormalTok{)))}
    \DecValTok{68}              \BuiltInTok{print}\NormalTok{(}\SpecialStringTok{f"N=}\SpecialCharTok{\{}\NormalTok{N}\SpecialCharTok{\}}\SpecialStringTok{: 既存ファイルあり（上書きせず検証のみ）一致=}\SpecialCharTok{\{}\NormalTok{same\_prev}\SpecialCharTok{\}}\SpecialStringTok{ residual=}\SpecialCharTok{\{}\NormalTok{residual}\SpecialCharTok{:.3e\}}\SpecialStringTok{ }\SpecialCharTok{\{}\NormalTok{status}\SpecialCharTok{\}}\SpecialStringTok{"}\NormalTok{)}
    \DecValTok{69}              \ControlFlowTok{if} \KeywordTok{not}\NormalTok{ same\_prev:}
    \DecValTok{70}\NormalTok{                  gate\_ok }\OperatorTok{=} \VariableTok{False}
    \DecValTok{71}\NormalTok{                  status }\OperatorTok{=} \StringTok{"EXISTING\_MISMATCH"}
    \DecValTok{72}          \ControlFlowTok{else}\NormalTok{:}
    \DecValTok{73}\NormalTok{              np.savez\_compressed(out\_path, v}\OperatorTok{=}\NormalTok{v, g}\OperatorTok{=}\NormalTok{g, Z0}\OperatorTok{=}\NormalTok{Z,}
    \DecValTok{74}\NormalTok{                                  sigma}\OperatorTok{=}\NormalTok{sig, residual}\OperatorTok{=}\NormalTok{np.float64(residual),}
    \DecValTok{75}\NormalTok{                                  n}\OperatorTok{=}\NormalTok{np.int64(N), seed}\OperatorTok{=}\NormalTok{np.int64(SEED),}
    \DecValTok{76}\NormalTok{                                  delta}\OperatorTok{=}\NormalTok{np.float64(DELTA), tol}\OperatorTok{=}\NormalTok{np.float64(TOL),}
    \DecValTok{77}\NormalTok{                                  iters}\OperatorTok{=}\NormalTok{np.int64(ITERS))}
\end{Highlighting}
\end{Shaded}

\begin{Shaded}
\begin{Highlighting}[]
    \DecValTok{81}      \ControlFlowTok{with} \BuiltInTok{open}\NormalTok{(os.path.join(PARENT\_DIR, }\StringTok{"parents\_summary.csv"}\NormalTok{), }\StringTok{"w"}\NormalTok{, newline}\OperatorTok{=}\StringTok{""}\NormalTok{) }\ImportTok{as}\NormalTok{ fh:}
    \DecValTok{82}\NormalTok{          w }\OperatorTok{=}\NormalTok{ csv.writer(fh)}
    \DecValTok{83}\NormalTok{          w.writerow([}\StringTok{"N"}\NormalTok{, }\StringTok{"M"}\NormalTok{, }\StringTok{"parent\_residual"}\NormalTok{, }\StringTok{"rank\_planes"}\NormalTok{, }\StringTok{"status"}\NormalTok{])}
    \DecValTok{84}\NormalTok{          w.writerows(rows)}
    \DecValTok{85}\NormalTok{      n\_bad }\OperatorTok{=} \BuiltInTok{sum}\NormalTok{(}\DecValTok{1} \ControlFlowTok{for}\NormalTok{ r }\KeywordTok{in}\NormalTok{ rows }\ControlFlowTok{if}\NormalTok{ r[}\DecValTok{4}\NormalTok{] }\OperatorTok{!=} \StringTok{"ok"}\NormalTok{)}
    \DecValTok{86}      \BuiltInTok{print}\NormalTok{(}\SpecialStringTok{f"summary: }\SpecialCharTok{\{}\BuiltInTok{len}\NormalTok{(rows)}\SpecialCharTok{\}}\SpecialStringTok{ parents, }\SpecialCharTok{\{}\NormalTok{n\_bad}\SpecialCharTok{\}}\SpecialStringTok{ non{-}ok"}\NormalTok{)}
    \DecValTok{87}      \ControlFlowTok{if} \KeywordTok{not}\NormalTok{ gate\_ok:}
    \DecValTok{88}          \BuiltInTok{print}\NormalTok{(}\StringTok{"GATE FAIL"}\NormalTok{)}
    \DecValTok{89}\NormalTok{          sys.exit(}\DecValTok{1}\NormalTok{)}
    \DecValTok{90}      \BuiltInTok{print}\NormalTok{(}\StringTok{"STATIC PARENTS DONE"}\NormalTok{)}
\end{Highlighting}
\end{Shaded}

\textbf{Exceptional behavior}: - Existing files are \textbf{never
overwritten}. If a file exists, only bit-identity verification against
the regenerated values is performed (lines 65--71); a mismatch is
flagged EXISTING\_MISMATCH and fails the gate. - Non-convergence
(residual ≥ 1e-8) is not an exception; it is explicitly recorded in the
ledger as NOT\_CONVERGED (lines 54--55). - Any gate failure yields exit
code 1 (lines 87--89).

\subsection{5. Execution Results}\label{execution-results}

\subsubsection{5.1 Reproduction Commands}\label{reproduction-commands}

\begin{Shaded}
\begin{Highlighting}[]
\BuiltInTok{cd}\NormalTok{ N3\_N40\_stage123\_sweep\_20260905}
\ExtensionTok{python3}\NormalTok{ make\_static\_parents\_N3\_N40\_v1.py      }\CommentTok{\# standalone}
\CommentTok{\# or as the first step of ./run\_all.sh}
\end{Highlighting}
\end{Shaded}

\subsubsection{5.2 Execution Environment}\label{execution-environment}

\begin{itemize}
\tightlist
\item
  Python 3.9.6 (\texttt{.venv/bin/python3})
\item
  numpy 2.0.2 (BLAS/LAPACK: macOS Accelerate)
\item
  macOS 26.3.1 (arm64, Darwin 25.x)
\item
  RNG: numpy \texttt{default\_rng} (PCG64)
\end{itemize}

\subsubsection{5.3 Execution Time}\label{execution-time}

\textbf{Roughly 30 seconds to 1 minute} for all 38 parents (including
generation, saving, verification, and ledger output). The dominant cost
per parent is the eigendecomposition of JG (2N×2N) times the iteration
count; \textasciitilde0.15 s per parent construction at N=40.

\subsubsection{5.4 Verification Gates}\label{verification-gates}

{\def\LTcaptype{none} % do not increment counter
\begin{longtable}[]{@{}
  >{\raggedright\arraybackslash}p{(\linewidth - 6\tabcolsep) * \real{0.2500}}
  >{\raggedright\arraybackslash}p{(\linewidth - 6\tabcolsep) * \real{0.2500}}
  >{\raggedright\arraybackslash}p{(\linewidth - 6\tabcolsep) * \real{0.2500}}
  >{\raggedright\arraybackslash}p{(\linewidth - 6\tabcolsep) * \real{0.2500}}@{}}
\toprule\noalign{}
\begin{minipage}[b]{\linewidth}\raggedright
Gate
\end{minipage} & \begin{minipage}[b]{\linewidth}\raggedright
Pass condition
\end{minipage} & \begin{minipage}[b]{\linewidth}\raggedright
Measured
\end{minipage} & \begin{minipage}[b]{\linewidth}\raggedright
Verdict
\end{minipage} \\
\midrule\noalign{}
\endhead
\bottomrule\noalign{}
\endlastfoot
G1: N=40 lineage & generated v, g, Z0 equal (via
\texttt{np.array\_equal}) to the canonical static parent (20260904;
itself verified bit-identical to the July canonical run) & equal &
\textbf{PASS} \\
G1': file identity & (observation) SHA256 of the N=40 npz equals that of
the canonical static-parent file & both \texttt{eadc87ee…}, identical &
\textbf{PASS} \\
G2: convergence & residual \textless{} 1e-8 for all N (status=ok) &
38/38 ok, residual ∈ {[}3.54e-14, 9.42e-13{]} & \textbf{PASS} \\
G3: execution & exit code 0 and \texttt{STATIC\ PARENTS\ DONE} printed &
confirmed & \textbf{PASS} \\
\end{longtable}
}

\subsubsection{5.5 Data}\label{data}

{\def\LTcaptype{none} % do not increment counter
\begin{longtable}[]{@{}
  >{\raggedright\arraybackslash}p{(\linewidth - 2\tabcolsep) * \real{0.5000}}
  >{\raggedright\arraybackslash}p{(\linewidth - 2\tabcolsep) * \real{0.5000}}@{}}
\toprule\noalign{}
\begin{minipage}[b]{\linewidth}\raggedright
Item
\end{minipage} & \begin{minipage}[b]{\linewidth}\raggedright
Content
\end{minipage} \\
\midrule\noalign{}
\endhead
\bottomrule\noalign{}
\endlastfoot
Folder & \texttt{N3\_N40\_stage123\_sweep\_20260905/parents/} \\
Static-parent npz &
\texttt{parent\_static\_N00003..N00040\_makeparent\_20260905.npz} (38
files) \\
Sizes & \textasciitilde2.0KB (N=3; 1,980 bytes) to \textasciitilde37KB
(N=40; 37,396 bytes), total \textasciitilde636KB \\
Ledger & \texttt{parents\_summary.csv} (38 rows: N, M, parent\_residual,
rank\_planes, status) \\
SHA256 & canonical values for all files in the bundled
\texttt{SHA256SUMS.txt}. Representative values: \\
\end{longtable}
}

\begin{verbatim}
c3230103e82976decde7bbe6fe5df545d12804127cf274f0d385e044ce544ca2  make_static_parents_N3_N40_v1.py
ba0fc19b03caf06d16e97e2cd5da499ed2ed8e95288f6dbf2062bedcaf11176d  run_n_scaling_lowrank_v1.py
eadc87ee0276554c7ab02e571e05200f0b719c1250b82607c7546b30a4d6f232  parents/parent_static_N00040_makeparent_20260905.npz
f384d9e90cc0a3a87a82a9a6928c9996d51bd0f83862c79f2cfe8a5ce182bbca  parents/parents_summary.csv
\end{verbatim}

\subsubsection{5.6 Figures}\label{figures}

No figures are generated in this chapter (the step-0 grid figure of
Chapter 3, \texttt{fig\_complex\_plane\_step0\_N3\_N40\_stage123.png},
visualizes this chapter's products).

\subsection{6. Execution Analysis (objective report and observations
only)}\label{execution-analysis-objective-report-and-observations-only}

\begin{enumerate}
\def\labelenumi{\arabic{enumi}.}
\tightlist
\item
  All 38 parents converged within the first restart; the eigenmode
  residuals fell in the range 3.54×10\^{}-14 to 9.42×10\^{}-13 (all at
  or below the order of tol=1e-12). Zero cases of NOT\_CONVERGED or
  GATE\_FAIL.
\item
  The N=40 products are bit-identical, at the array level, to the
  canonical static parent generated independently (on 2026-09-04, in a
  different folder), and the SHA256 of the npz files themselves also
  coincide. This is direct evidence that the generation is fully
  deterministic under the same seed, engine, and environment.
\item
  The number of nonzero planes of the parent's σ spectrum (rank\_planes)
  was observed as N=3: 2, N=4: 2, N=5: 4, N=6: 6, and N for N≥7 (see the
  ledger \texttt{parents\_summary.csv}).
\item
  File sizes grow roughly in proportion to M = N(N−1)/2.
\item
  The rng consumption is the minimal ``m uniforms + 2m normals'' (when
  no restart occurs); every run completed with this minimal consumption
  (the bit-identity of item 2 is corroborating evidence).
\end{enumerate}

\begin{center}\rule{0.5\linewidth}{0.5pt}\end{center}

(End of Chapter 1. Chapter 2, ``The Sweep --- Stage 1+2+3 Dynamics'',
follows.)

\end{document}
